\chapter{State of the art}
\label{ch:sota}

Predicting where a person goes next is an old question with a large literature.
Eight published works matter to this report, and this chapter says, in plain
terms, what each one gave me. Two of them set the frame every later result is
read in, three place the instrument, two are the closest relatives of my one
substantial gain, and one marks the boundary between the published task and the
task I actually took on.

\section{How predictable is a person?}
\label{sec:sota-ceiling}

The most useful part of this literature is not the part that proposes
predictors. It is the part that measures how much room a predictor has before it
even starts.

Song et al.\ measured exactly that \cite{song2010limits}. They took three months
of mobile-phone records for 45{,}000 people, treated each person's sequence of
visited antennas as a stream of symbols, and asked how much of the next position
is already implied by the order of everything that came before. From that
quantity they derived a \emph{ceiling}: the best score any predictor could ever
reach on that person, however clever it is. The ceiling came out narrowly peaked
at 93\,\%.

Three things in that paper matter here, and the 93\,\% is the least of them.

\textbf{The ceiling is the same for everybody.} The authors say plainly that they
expected heavy travellers to be harder to predict than people who never leave
their neighbourhood, and they were wrong: past a radius of ten kilometres, the
ceiling stops depending on distance travelled at all. Whatever separates an easy
person from a hard one, it is not how much ground they cover.

\textbf{The same ceiling does not mean the same score.} The same paper also
computes what can be predicted from a person's visit frequencies alone, ignoring
the order in which places were visited. That quantity is spread very widely
across the population and peaks near 30\,\%. So the room is the same for
everyone, while what a predictor manages to extract from it is not. That is
exactly the quantity Chapters~\ref{ch:user} and~\ref{ch:groups} measure on my
own data, where it runs from 0 to 100\,\% depending on the person.

\textbf{And a curve published almost in passing, which cost me an iteration.}
Splitting the week into 168 hourly slots, the authors measure how often a person
is found in their most-visited place for that hour of the week. The answer is
about 70\,\% on average, but close to 90\,\% during the night, when people are
reliably at home, with clear low points around midday and between six and seven
in the evening. In other words, the hours at which people are hard to predict
were published in 2010, and they are not the small hours of the morning.
Chapter~\ref{ch:time} spends its whole length rediscovering that the expensive
way.

\section{Most of that predictability is just standing still}

A ceiling of 93\,\% sounds like an invitation. The most useful thing I read
during the internship is the argument that it is not one.

Ikanovic and Mollgaard make that argument precisely
\cite{ikanovic2017alternative}. The quantity Song et al.\ bound answers the
question \emph{``where will this person be during the next time slot''}. In most
time slots the answer is ``exactly where they are now'', because people do not
move most of the time. So most of the apparent 93\,\% is not knowledge about
anybody's habits; it is the fact that staying put is the normal case. The
authors re-ran the same machinery on the question a predictor is actually useful
for, \emph{``where will this person go next''}, by removing the repetitions
before measuring. On location traces from 604 people, the ceiling fell from
above 90\,\% to \textbf{71\,\%}. Roughly a quarter of what looked like
predictability was the passage of time rather than the behaviour of a person.

This is what names my competitor, and it is why the whole report is written
against one. If most of a high score is people staying where they are, then the
rule ``answer the last antenna seen'' already collects most of it, for free. On
an ordinary day of my corpus that rule is right \textbf{69.2\,\%} of the time
(Chapter~\ref{ch:data}), which is roughly where Song et al.'s hourly regularity
sits, and well below the 93\,\% ceiling. The distance between those two numbers
is the entire working space of this report.

Two conventions follow from that, and they are used everywhere after this page.

\paragraph{Every result is a distance, never an absolute score.} The same trivial
rule scores 69.2\,\% on one of the files used here and 43.1\,\% on another,
because the people in them differ (Table~\ref{tab:datasets}). An absolute
accuracy therefore describes the test group more than it describes the model.
Only the distance to the rule, measured on the same moments, survives that.

\paragraph{The interesting quantity is per kind of person, not on average.} If
the ceiling is the same for everybody but the score is not, a population average
hides the thing worth reporting. Chapters~\ref{ch:user} and~\ref{ch:groups}
measure the spread directly: 22.4 correct answers per hundred separate the four
kinds of user, against 0.4 across an entire range of training sizes. That
decomposition is this report's main claim.

\section{What other people have used to predict}

The predictors themselves I take from one place. Luca et al.\ survey the whole
deep-learning branch of human mobility, from networks that read a trajectory
step by step\footnote{\textbf{Recurrent network}: a network that reads a sequence
one item at a time while carrying a summary of what it has already read.} to
mechanisms that let a prediction reach back to a matching moment in an earlier
day \cite{luca2021survey}. Three points in that survey bear on the chapters
below.

\textbf{Giving a model the time of day is standard practice}, and has been for
years, so the mechanism of Chapter~\ref{ch:time} is not an idea of mine. What
Chapters~\ref{ch:time} and~\ref{ch:groups} add is the finding that \emph{which}
hours matter is a property of the person rather than of the corpus.

\textbf{Scoring a prediction at several depths is the convention of the field},
and it is not something I invented either. Predictors in this literature are
reported at top-1, top-3 and top-5: whether the correct answer is the model's
first guess, or somewhere in its first three or first five \cite{luca2021survey}.
The language-model work that is closest to mine reports the same three depths
\cite{qin2025mobllm}, which is why this report does too, and
Chapter~\ref{ch:first} shows the three depths on one picture. The reason the
convention exists is practical: an operator that can prepare three antennas
instead of one is served by a correct answer in third place, so a model that
never puts the right antenna first but often puts it third is still worth
something.

\textbf{And the reaching-back mechanism} that the better step-by-step networks
had to be extended with is the native operation of the architecture behind
current language models,\footnote{\textbf{Transformer}: the architecture behind
current language models, in which every position of a sequence can look directly
at every earlier one.} which is one honest reason to reach for an existing
language model here instead of building a predictor from scratch.

\section{Two ways to use a language model}

Since 2023 the field has been reaching for those models, and it split into the
same two branches the team I joined was split into (Chapter~\ref{ch:intro}).

The first branch leaves the model's numbers alone and writes a better
instruction. Wang et al.\ turn a person's recent stops into a paragraph of
English, separating the long-term history from the immediate context, and ask a
large commercial model for the next location \cite{wang2023llmmob}. They get
results comparable with purpose-built predictors without training anything, and
they report that the design of the instruction, not the choice of model, is what
moves the score. That is in substance the branch Tiphain and Benjamin were
taking on this corpus.

The second branch changes the numbers, and it is mine. Qin et al.\ fine-tune a
small open-source model with the same cheap method used here, across six
mobility datasets, and argue that a model adapted this way transfers to a new
corpus better than prompting a large commercial one does, at a fraction of the
running cost \cite{qin2025mobllm}. It is the closest published relative of what
this report attempts.

\section{Sorting a population into kinds of user}

The step that produced this report's only substantial gain, sorting people by
\emph{when} they move and letting each kind inform its own members
(Chapter~\ref{ch:groups}), has two ancestors I found only after the fact.

Pappalardo et al.\ establish the observation \cite{pappalardo2015returners}.
Looking at both mobile-phone and location data, they find that a population does
not spread smoothly along a single scale of mobility. It separates into two
recognisable kinds of person: \emph{returners}, who keep cycling between a small
set of habitual places however far they travel, and \emph{explorers}, whose set
of visited places keeps growing. The two behave differently enough that a
phenomenon spreading through the population moves at a different speed depending
on which kind carries it. The four kinds of user in Chapter~\ref{ch:groups} are a
version of the same finding with the axis changed: what separates mine is not
how far a person travels but at which hours they change location.

Jeong, Leconte and Proutière supply the mechanism \cite{jeong2016cluster}. Their
argument is that any single person's history is too short to estimate that
person's own habits from, so they group people who resemble each other and
predict one person's next location partly from the trajectories of the others.
The property that makes it more than a heuristic is that their method falls back
to per-person prediction when the resemblances are not really there, so grouping
cannot make things worse. Chapter~\ref{ch:groups} does the same thing in a
different place, and spends most of its length on one constraint they also face:
the group has to be recoverable from the person's past alone.

\section{Predicting an antenna rather than a place}

Everything above predicts a \emph{place}: a home, a workplace, a stopping point
obtained by grouping nearby records together. My target is the next antenna,
which is a different object, and the work that predicts antennas rather than
places sits in the networking community and asks the question for another
purpose. Chen, Mériaux and Valentin predict which antenna a phone will be handed
over to,\footnote{\textbf{Handover}: the moment the network moves an ongoing
connection from one antenna to another, normally because the new one now offers
a better signal.} from radio measurements, so that the network can prepare the
handover before it happens \cite{chen2013nextcell}. They have signal strengths my
corpus does not carry, and they do not need to know where the person is going,
only which antenna will serve them next.

Predicting an antenna from a trace of antennas, with no radio measurements, sits
between the two and inherits a drawback from each. In this corpus, 12.4\,\% of
back-to-back records change antenna without the person having really gone
anywhere, because the phone re-attached to another antenna of the same mast
(Chapter~\ref{ch:data}), so part of what my model is asked to predict is a
decision of the network rather than of a human being. And the human part of it is
graded on antennas rather than on places, which is a finer target than the
published ceilings were ever estimated for.
